\chapter*{Abstract}
\addcontentsline{toc}{chapter}{Abstract}

Researchers increasingly need to locate, read and compare scientific papers faster than manual reading allows,
yet general-purpose language models answer questions without showing where an answer came from and may state
unsupported facts. This report describes the design, implementation and verification of the
\emph{Scientific Literature Intelligence Platform}, a single-user web application that answers natural-language
questions over a corpus of scientific papers and cites every claim to a specific passage and page.

Papers are imported from arXiv or uploaded as PDF files. The system extracts text page by page with PyMuPDF,
splits it into 256-token chunks that keep page and character provenance, embeds the chunks with the
\code{all-MiniLM-L6-v2} sentence encoder and stores them in PostgreSQL with the pgvector extension. Retrieval
combines dense vector search with BM25-style full-text search through reciprocal rank fusion and re-orders the
candidates with a cross-encoder. Answers are generated by a locally hosted \code{qwen2.5:3b} model served by
Ollama; the server resolves every citation against the passages actually supplied to the model and returns an
explicit ``insufficient evidence'' response when the evidence does not support an answer. The system also
produces single-paper summaries, cross-paper syntheses, structured extractions and comparison tables. It is
implemented with FastAPI, SQLAlchemy and Alembic, a React and TypeScript frontend, and Docker Compose, and is
checked by a GitHub Actions pipeline.

On a version-controlled evaluation set of 50 questions over 20 natural-language-processing papers, the default
retrieval configuration reached a Recall@5 of 0.875 and a mean reciprocal rank of 0.701; every citation returned
in two complete question-answering runs referred to evidence supplied to the model, and nine of ten unanswerable
questions were declined. The 95th-percentile question-answering latency on the reference laptop was 4.4--4.5\,s,
above the 3\,s target. The evaluation labels were written by an AI assistant and have not yet been reviewed by a
person, the retrieval configuration was selected on the same questions, and spot checks found wrong-paper
answers and claims not supported by their cited passage. These limitations are reported alongside the results.

\vspace{1em}
\noindent\textbf{Keywords:} retrieval-augmented generation, semantic search, citation grounding, local language
models, scientific literature, software engineering.
